% Auto-generated from ch01-basics.mdx — edit ch01-basics.mdx then re-run mdx2tex.mjs
\chapter{第 1 章 装饰器基础}
\label{ch:ch01-basics}

\begin{tcolorbox}[colback=blue!5,colframe=blue!40,title={📘 本章地图}]
本章解决三个问题：为什么装饰器能成立（一等函数 + 闭包）？\textbackslash\{\}textbackslash\textbackslash\{\}\{\textbackslash\{\}\}texttt\textbackslash\{\}\{@\textbackslash\{\}\} 到底是什么语法？如何写第一个装饰器？

\begin{itemize}
\item 1.1 函数是一等公民 - 1.2 闭包：内部函数记住环境 - 1.3 \textbackslash\{\}textbackslash\textbackslash\{\}\{\textbackslash\{\}\}textbackslash\textbackslash\{\}textbackslash\textbackslash\{\}\{\textbackslash\{\}\}\textbackslash\{\}\{\textbackslash\{\}textbackslash\textbackslash\{\}\{\textbackslash\{\}\}\textbackslash\{\}\}texttt\textbackslash\{\}textbackslash\textbackslash\{\}\{\textbackslash\{\}\}\textbackslash\{\}\{@\textbackslash\{\}textbackslash\textbackslash\{\}\{\textbackslash\{\}\}\textbackslash\{\}\} 语法糖：\textbackslash\{\}textbackslash\textbackslash\{\}\{\textbackslash\{\}\}textbackslash\textbackslash\{\}textbackslash\textbackslash\{\}\{\textbackslash\{\}\}\textbackslash\{\}\{\textbackslash\{\}textbackslash\textbackslash\{\}\{\textbackslash\{\}\}\textbackslash\{\}\}texttt\textbackslash\{\}textbackslash\textbackslash\{\}\{\textbackslash\{\}\}\textbackslash\{\}\{f = deco(f)\textbackslash\{\}textbackslash\textbackslash\{\}\{\textbackslash\{\}\}\textbackslash\{\}\} - 1.4 Recipe：计时器装饰器 - 1.5 常见陷阱 - 1.6 要点总结
\end{itemize}
\end{tcolorbox}
\section{1.1 函数是一等公民}
"一等公民"（first-class citizen）意味着函数可以：

\begin{enumerate}
\item 赋值给变量
\item 作为参数传给另一个函数
\item 作为另一个函数的返回值
\end{enumerate}

\begin{codeblock}{python}
# file: examples/first_class.py
def greet(name: str) -> str:
    """返回一条问候语。"""
    return f"Hello, {name}!"

# 1) 赋值给变量
say_hello = greet
print(say_hello("Alice"))  # Hello, Alice!

# 2) 作为参数传入高阶函数
def shout(fn, text: str) -> str:
    """调用 fn(text) 并把结果转为大写。"""
    return fn(text).upper()

print(shout(greet, "bob"))  # HELLO, BOB!
\end{codeblock}

这三条性质是所有装饰器的"地基"。没有一等函数，装饰器根本无从谈起。

\section{1.2 闭包：内部函数记住环境}
\textbackslash\{\}textbf\{闭包\}（closure）是指：一个内部函数引用了外部函数的局部变量，即使外部函数已经返回，内部函数仍然能访问这些变量。


\begin{codeblock}{python}
# file: examples/closure.py
def make_counter() -> callable:
    """返回一个闭包：每次调用计数 +1。"""
    count = 0  # 这个变量被闭包"记住"

    def increment() -> int:
        # nonlocal 告诉 Python：count 来自外层作用域
        nonlocal count
        count += 1
        return count

    return increment

# 调用：counter 就是 increment，却记住了 count
counter = make_counter()
print(counter())  # 1
print(counter())  # 2
print(counter())  # 3
\end{codeblock}

关键点：\textbackslash\{\}texttt\{counter()\} 每次都能读到自己那份 \textbackslash\{\}texttt\{count\}，即使 \textbackslash\{\}texttt\{make\_counter\} 早已执行完毕。装饰器就是"带额外行为的闭包"。

\section{1.3 \textbackslash\{\}texttt\{@\} 语法糖：\textbackslash\{\}texttt\{f = deco(f)\}}
PEP 318 引入的 \textbackslash\{\}texttt\{@\} 语法只是\textbackslash\{\}textbf\{语法糖\}（syntactic sugar）。下面两种写法\textbackslash\{\}textbf\{完全等价\}：


\begin{codeblock}{python}
# 写法 A：语法糖
@log
def hello(name):
    return f"hi {name}"

# 写法 B：手动调用
def hello(name):
    return f"hi {name}"
hello = log(hello)
\end{codeblock}

也就是说，装饰器是一个\textbackslash\{\}textbf\{输入为函数、输出为另一个函数\}的高阶函数：

\textbackslash\{\}[ T: f \textbackslash\{\}mapsto f' \textbackslash\{\}]

其中 \$f'\$ 通常是一个"包裹了 \$f\$ 的 wrapper"。

\subsection{最小可用装饰器}

\begin{codeblock}{python}
# file: examples/log_deco.py
import functools

def log(func):
    """打印函数调用日志的装饰器。"""

    # functools.wraps 拷贝原函数的 __name__、__doc__ 等元信息
    @functools.wraps(func)
    def wrapper(*args, **kwargs):
        print(f"[log] 调用 {func.__name__}, args={args}, kwargs={kwargs}")
        result = func(*args, **kwargs)  # 调用原函数
        print(f"[log] {func.__name__} 返回 {result!r}")
        return result

    return wrapper

@log
def add(a: int, b: int) -> int:
    return a + b

add(2, 3)
# [log] 调用 add, args=(2, 3), kwargs={}
# [log] add 返回 5
\end{codeblock}

图 1-1 是 \textbackslash\{\}texttt\{@log\} 装饰器的调用时序：

\textbackslash\{\}begin\{figure\}[htbp]\textbackslash\{\}centering\textbackslash\{\}includegraphics[width=0.85\textbackslash\{\}textwidth]\{figures/decorator-sequence.svg\}\textbackslash\{\}caption\{图 1-1：@log 装饰器调用时序\}\textbackslash\{\}end\{figure\}

如图所示，每次调用 \textbackslash\{\}texttt\{add(2, 3)\} 实际上进入的是 \textbackslash\{\}texttt\{wrapper\}，它先打日志、再调原函数、最后返回结果。

\subsection{装饰器 vs 普通函数}
\begin{table}[htbp]
\centering
\begin{tabularx}{\textwidth}{X X X}
\toprule
维度 & 普通函数 & 装饰器 \\
\midrule
输入 & 数据 & \textbackslash\{\}textbf\{另一个函数\} \\
输出 & 数据 & \textbackslash\{\}textbf\{包裹后的函数\} \\
用法 & \textbackslash\{\}texttt\{result = f(x)\} & \textbackslash\{\}texttt\{f = deco(f)\}（或 \textbackslash\{\}texttt\{@deco\}） \\
典型场景 & 业务逻辑 & 横切关注点：日志/缓存/权限 \\
\bottomrule
\end{tabularx}
\end{table}

\section{1.4 Recipe：计时器装饰器}
<div className="callout recipe">

\subsubsection{🍳 Recipe 1-1：测量函数耗时}
\textbackslash\{\}textbf\{Problem\}：想知道一个函数运行了多久，又不想把计时代码散落到每处调用。

\textbackslash\{\}textbf\{Solution\}：


\begin{codeblock}{python}
# file: recipes/timer.py
import time
import functools

def timer(func):
    @functools.wraps(func)
    def wrapper(*args, **kwargs):
        t0 = time.perf_counter()       # 开始计时（高精度）
        try:
            return func(*args, **kwargs)
        finally:
            dt = time.perf_counter() - t0
            print(f"[timer] {func.__name__} 耗时 {dt*1000:.2f} ms")
    return wrapper

@timer
def slow_sum(n: int) -> int:
    return sum(range(n))

slow_sum(10_000_000)  # [timer] slow_sum 耗时 253.47 ms
\end{codeblock}

\textbackslash\{\}textbf\{Discussion\}：

\begin{itemize}
\item 用 \textbackslash\{\}texttt\{try/finally\} 保证即使函数抛异常也会打印耗时。
\item \textbackslash\{\}texttt\{time.perf\_counter()\} 是 Python 推荐的短时计时 API（比 \textbackslash\{\}texttt\{time.time()\} 精度更高，且不受系统时钟回拨影响）。
\end{itemize}
\textbackslash\{\}textbf\{See Also\}：第 2 章将把它升级为带参数的 \textbackslash\{\}texttt\{@timer(unit="ms")\}。

</div>

\section{1.5 常见陷阱}
<div className="callout pitfall">

\subsubsection{⚠️ 陷阱 1：忘记 \textbackslash\{\}texttt\{@functools.wraps(func)\}}
如果不加 \textbackslash\{\}texttt\{@functools.wraps(func)\}，装饰后的函数会丢失原函数的元信息：


\begin{codeblock}{python}
@log
def add(a, b):
    """两个整数相加。"""
    return a + b

print(add.__name__)  # 没有 wraps → "wrapper"，加了 wraps → "add"
print(add.__doc__)   # 没有 wraps → None，加了 wraps → "两个整数相加。"
\end{codeblock}

这会破坏调试器、文档生成工具（Sphinx）、以及依赖 \textbackslash\{\}texttt\{\_\_name\_\_\} 的框架（如 Flask 路由注册）。

\textbackslash\{\}textbf\{正确做法\}：\textbackslash\{\}textbf\{永远\}在 wrapper 上加 \textbackslash\{\}texttt\{@functools.wraps(func)\}。

</div>

\section{1.6 要点总结}
\begin{tcolorbox}[colback=green!5,colframe=green!40,title={✅ 核心要点}]
\begin{itemize}
\item 装饰器的前提是\textbackslash\{\}textbackslash\textbackslash\{\}\{\textbackslash\{\}\}textbf\textbackslash\{\}\{函数是一等公民\textbackslash\{\}\} + \textbackslash\{\}textbackslash\textbackslash\{\}\{\textbackslash\{\}\}textbf\textbackslash\{\}\{闭包\textbackslash\{\}\}。 - \textbackslash\{\}textbackslash\textbackslash\{\}\{\textbackslash\{\}\}texttt\textbackslash\{\}\{@deco\textbackslash\{\}\} 只是 \textbackslash\{\}textbackslash\textbackslash\{\}\{\textbackslash\{\}\}texttt\textbackslash\{\}\{f = deco(f)\textbackslash\{\}\} 的语法糖。 - 一个合格的装饰器必须用 \textbackslash\{\}textbackslash\textbackslash\{\}\{\textbackslash\{\}\}texttt\textbackslash\{\}\{@functools.wraps(func)\textbackslash\{\}\} 保留原函数元信息。 - wrapper 用 \textbackslash\{\}textbackslash\textbackslash\{\}\{\textbackslash\{\}\}texttt\textbackslash\{\}\{\textbackslash\{\}textbackslash\textbackslash\{\}\{\textbackslash\{\}\}emph\textbackslash\{\}\{args, \textbackslash\{\}textbackslash\textbackslash\{\}\{\textbackslash\{\}\}textbf\textbackslash\{\}\{kwargs\textbackslash\{\}\} 才能适配任意签名的被装饰函数。 - 装饰器是\textbackslash\{\}\}横切关注点\textbackslash\{\}\}*（日志、计时、权限、缓存）的首选表达。
\end{itemize}
\end{tcolorbox}
\section{1.7 延伸阅读}
\begin{itemize}
\item 第 2 章 · 进阶装饰器：带参装饰器、类装饰器、堆叠顺序
\item 附录 A · functools 速查
\item PEP 318 – Decorators for Functions and Methods：https://peps.python.org/pep-0318/
\end{itemize}
---

\textbackslash\{\}textbf\{本章参考文献\}

\begin{enumerate}
\item Python Software Foundation. "PEP 318 – Decorators for Functions and Methods." https://peps.python.org/pep-0318/
\item Real Python. "Primer on Python Decorators." https://realpython.com/primer-on-python-decorators/
\end{enumerate}
